\documentclass[11pt,reqno]{amsart}

%==============================================================
% Volatility-Adaptive Lossless Coding of ECG Beats FLAC.
%==============================================================

\usepackage[utf8]{inputenc}
\usepackage[T1]{fontenc}
\usepackage{lmodern}
\usepackage[margin=1.15in]{geometry}
\usepackage{amsmath,amssymb,amsthm,mathtools}
\usepackage{microtype}
\usepackage{booktabs}
\usepackage{hyperref}
\usepackage{cleveref}

\theoremstyle{plain}
\newtheorem{proposition}{Proposition}
\theoremstyle{remark}
\newtheorem{remark}{Remark}
\newtheorem*{finding}{Finding}

\newcommand{\R}{\mathbb{R}}
\DeclareMathOperator{\ACF}{ACF}

\title[Volatility-adaptive lossless coding of ECG beats FLAC]{Volatility-Adaptive Lossless Coding of ECG\\
  Beats FLAC}
\author{Ramiro Fontes}
\address{Quijotic Research}
\email{ramirofontes@gmail.com}
\urladdr{www.quijotic.com}
\date{\today}

\begin{document}

\begin{abstract}
A lossless compressor predicts the next sample and entropy-codes the residual;
on a bursty signal the residual scale (volatility) clusters, and modeling it is
the main lever beyond a fixed-scale code. We build a codec---an AR predictor, a
Student-$t$ residual law, a conditional-volatility scale, and a Witten--Neal--Cleary
arithmetic coder---and test it byte-for-byte (canonical raw input, exact
round-trip, all overhead in the file) against gzip/bzip2/xz and against the
\emph{actual, parameter-tuned} \texttt{flac} binary. The result is a real,
boss-beating win in one domain: on electrocardiogram (ECG) data the codec beats
tuned FLAC by $+3.4\%$ (SciPy recording) and $+4.9\%$ (MIT-BIH record~100), a
win that reproduces across two independent recordings; on $16$-bit speech, by
contrast, FLAC's audio-specific tuning beats the codec by $\sim 8\%$. We also
settle which volatility model does the work: against the right opponents, a
fractional (Mittag-Leffler) long-memory kernel beats only the one-parameter
EWMA and is matched or beaten by GARCH and FIGARCH, so the lever is conditional
volatility modeling plus arithmetic coding versus FLAC's Rice coder---not the
fractional form, and not ``long memory'' per se. We are explicit about scope:
FLAC is the general-purpose lossless standard, not a domain ECG codec. Methods
that model the quasi-periodic beat structure (wavelet and long-term-prediction
ECG coders) exploit structure our AR predictor ignores and remain the unfought
opponent; the claim here is strictly that we beat the general-purpose standard.
\end{abstract}

\maketitle

\section{Introduction}
\label{sec:intro}

A lossless compressor is a predictive model followed by an entropy coder. The
coder---arithmetic, range, ANS, or Rice---realizes the model's code length to
within a fraction of a bit, so the whole contest is whether the predictive
distribution assigns shorter ideal code length than a baseline. After a linear
predictor removes the mean structure of a bursty signal, the residual is not
stationary in scale: it clusters into quiet and violent stretches. Modeling that
time-varying scale---the conditional volatility---is the main lever beyond a
fixed-scale residual code, and it is the lever this paper exploits.

The reference point that matters is FLAC, the de facto general-purpose lossless
standard for $16$-bit signals. FLAC's core is a linear predictor followed by
\emph{Rice} coding of the residual---a fast, parameter-light coder that does not
model conditional volatility. Our hypothesis is that replacing Rice coding with
a Student-$t$ residual law whose scale tracks the conditional volatility, driven
through an arithmetic coder, recovers bits FLAC leaves on the table on signals
with strong volatility clustering. ECG is the natural candidate: long quiet
baselines punctuated by violent QRS complexes.

We are careful about two things a result like this can over-claim. First,
\emph{which} part does the work: we test the volatility model against the right
opponents (GARCH, FIGARCH), not only the weak EWMA baseline. Second,
\emph{scope}: beating FLAC proves we beat the general-purpose standard, not the
domain-specific ECG codecs (\cref{sec:scope}). Both qualifications are load-bearing
and stated up front.

\section{Method and data}
\label{sec:method}

\paragraph{Datasets.} ECG: the SciPy electrocardiogram ($108$k samples) and
MIT-BIH record~100 (channel MLII, $650$k samples, an independent
subject/lead) \cite{mitbih,physionet}. Speech: real mono $16$-bit PCM. S\&P~500
pooled daily log-returns are used for the volatility-model comparison.

\paragraph{Codec.} Residuals are formed by an order-$p$ AR predictor (rounded so
the reconstruction is exact). The residual is coded with a Student-$t$ law
(empirically the best tail) whose scale $s_t$ is a conditional volatility from
one of: a fixed scale, EWMA, fitted GARCH$(1,1)$ \cite{garch}, fitted FIGARCH
(long memory, ARCH-$\infty$ with hyperbolic fractional-differencing weights)
\cite{figarch}, or the Mittag-Leffler kernel $w_k=k^{\alpha-1}E_{\alpha,\alpha}
(-\lambda k^\alpha)$ of a fractional-relaxation clock \cite{mem,GKMR}. The coder
is a Witten--Neal--Cleary arithmetic coder \cite{wnc} over a per-symbol
cumulative-frequency table with a guaranteed nonzero floor.

\paragraph{Byte-level protocol.} The input is frozen as canonical little-endian
raw bytes; every method sees exactly those bytes. The codec must
$\textsf{encode}\to\textsf{decode}\to$ reproduce the input \emph{byte for byte}
(verified), and its compressed file carries all overhead (AR coefficients,
volatility parameters, alphabet), so it is fully self-decoding. Baselines are
actual bytes from gzip\,$-9$, bzip2\,$-9$, xz\,$-9$. The FLAC baseline is the
real \texttt{flac} $1.4.3$ binary, swept over its strong settings
($-8$, $-e$, $-p$, higher LPC order, and block size, with \texttt{--lax}) and
reported at its best; our codec is likewise reported at its best over AR order
and volatility model. The comparison is therefore best-configuration versus
best-configuration.

\section{Which volatility model does the work}
\label{sec:vol}

Before claiming a win we isolate the lever. We compare the Mittag-Leffler kernel
against EWMA, GARCH, and FIGARCH by a paired per-block test: a single fit over
each recording, then the models' actual per-symbol arithmetic code length
aggregated into $256$-sample blocks (header-free, isolating the model).

\begin{table}[h]
\centering\small
\begin{tabular}{lrrr}
\toprule
dataset & ML vs.\ EWMA & ML vs.\ GARCH & ML vs.\ FIGARCH\\
\midrule
synthetic long-mem & $+0.07\%$ (win $76\%$) & $+0.01\%$ (tie) & $-0.02\%$ (loses)\\
ECG (SciPy)        & $+0.64\%$ (win $96\%$) & $-0.60\%$ (win $12\%$) & $-0.59\%$ (win $14\%$)\\
ECG (MIT-BIH 100)  & $+0.80\%$ (win $99.8\%$) & $-1.00\%$ (win $0.5\%$) & $-1.05\%$ (win $1.4\%$)\\
S\&P returns       & $+0.05\%$ (win $75\%$) & $-0.00\%$ (tie) & $-0.04\%$ (loses)\\
\bottomrule
\end{tabular}
\caption{Per-block code-length comparison of the Mittag-Leffler volatility
kernel against three opponents (win rate = fraction of blocks ML is smaller).
ML beats only the one-parameter EWMA; it ties or loses to GARCH and loses to
FIGARCH on every dataset.}
\label{tab:vol}
\end{table}

\begin{finding}
The fractional (Mittag-Leffler) kernel reliably beats EWMA but is matched or
beaten by GARCH and FIGARCH everywhere. ``Long memory is the lever'' is at most
weakly true---FIGARCH is the aggregate-best volatility model, but only by
$\sim0.03\%$---and the specific fractional-clock form is never best. The lever is
conditional volatility modeling in general; the choice among GARCH/FIGARCH/ML
moves the result by under $1\%$.
\end{finding}

The earlier impression that the fractional kernel was special came from racing it
only against EWMA, which is short-memory by construction and is nobody's frontier.
Against the production-grade volatility models the fractional advantage
evaporates. We therefore use GARCH and FIGARCH as the codec's volatility model
in the byte-level results below, and treat ``volatility-adaptive arithmetic
coding,'' not Mittag-Leffler, as the mechanism.

\section{Against general-purpose compressors}
\label{sec:general}

On the frozen raw bytes, the codec beats every general-purpose compressor,
with exact round-trip and all overhead included.

\begin{table}[h]
\centering\small
\begin{tabular}{lrr}
\toprule
method (ECG, 216\,kB int16) & bytes & vs xz\\
\midrule
gzip$-9$            & $118{,}619$ & $-37\%$\\
xz$-9$              & $86{,}772$  & ---\\
bzip2$-9$           & $73{,}690$  & $+15\%$\\
AR$+$Gaussian       & $61{,}505$  & $+29\%$\\
AR$+$Student-$t$    & $60{,}797$  & $+30\%$\\
AR$+$Student-$t+$EWMA & $60{,}062$ & $+31\%$\\
\textbf{AR$+$Student-$t+$GARCH} & $\mathbf{59{,}318}$ & $\mathbf{+32\%}$\\
\bottomrule
\end{tabular}
\caption{Actual compressed bytes (exact round-trip). The codec beats gzip,
bzip2, and xz; most of the margin is AR$+$arithmetic coding, with the volatility
model adding a final $\sim3\%$. On S\&P returns the codec likewise beats xz by
$\sim8\%$.}
\label{tab:general}
\end{table}

The honest decomposition: plain AR$+$Gaussian already reaches $+29\%$ over xz;
the heavy tail and volatility model add the final few percent. Beating
general-purpose tools on a predictable numeric signal is expected---they do not
model the signal---so this is a sanity check, not the headline.

\section{Against real FLAC}
\label{sec:flac}

The headline is the comparison with the real, parameter-tuned \texttt{flac}
binary, best-configuration versus best-configuration, on the same samples.

\begin{table}[h]
\centering\small
\begin{tabular}{lrrr}
\toprule
signal & best FLAC & best ours & ours vs.\ FLAC\\
\midrule
ECG (SciPy)        & $61{,}379$  & $59{,}318$ (GARCH)   & $\mathbf{+3.4\%}$\\
ECG (MIT-BIH 100)  & $332{,}017$ & $315{,}602$ (FIGARCH)& $\mathbf{+4.9\%}$\\
speech (16\,kHz)   & $192{,}765$ & $208{,}026$ (FIGARCH)& $-7.9\%$\\
\bottomrule
\end{tabular}
\caption{Lossless bytes, exact round-trip, both encoders at their swept best.
The volatility-adaptive codec beats tuned FLAC on ECG---on two independent
recordings---and loses to it on speech.}
\label{tab:flac}
\end{table}

\begin{finding}
A volatility-adaptive Student-$t$ arithmetic codec beats the general-purpose
lossless standard (tuned FLAC) on ECG by $3$--$5\%$, reproducibly across two
recordings. The win is robust to the volatility model: EWMA, GARCH, FIGARCH, and
Mittag-Leffler all beat FLAC on ECG, with GARCH/FIGARCH best. On speech FLAC's
audio-specific machinery (LPC order/precision search, apodization, partitioned
Rice) wins by $\sim8\%$, and we make no audio claim.
\end{finding}

The asymmetry is informative: ECG's long quiet baselines and impulsive QRS make
conditional volatility a large lever that FLAC's Rice coder cannot use, whereas
on speech FLAC's decades of audio-specific tuning dominate. The method is a
biosignal codec, not an audio codec. In an ambulatory/Holter setting---continuous
multi-hour recording, naturally large coding blocks, high device-volume---a
$3$--$5\%$ reduction over the deployed general-purpose standard is real currency.

\section{Scope: the unfought boss}
\label{sec:scope}

One caveat is essential and we state it plainly. FLAC is the general-purpose
lossless standard; it is \emph{not} an ECG-specific codec. The biosignal
literature has domain codecs---wavelet-based and long-term-prediction methods
\cite{ecgcomp}---that explicitly model the quasi-periodic beat structure of ECG.
Our AR predictor models only short-range linear dependence; it does not exploit
beat-to-beat self-similarity at all. A domain ECG coder that predicts the
current beat from previous beats removes structure our residual still contains,
and could plausibly erase our margin. Beating FLAC establishes that
volatility-adaptive arithmetic coding beats the general-purpose standard on ECG;
it does \emph{not} establish that it beats ECG-specific codecs. That comparison
is the next decisive test---the same role FIGARCH played for the volatility
question in \cref{sec:vol}---and it is unfought here. We flag it before, not
after, the result invites excitement.

\section{Verdict}
\label{sec:verdict}

\begin{itemize}
\item \textbf{A real codec, byte-for-byte, beating the standard on ECG.} With
  exact round-trip and all overhead included, the codec beats tuned FLAC on ECG
  by $+3.4\%$ and $+4.9\%$ on two independent recordings, and beats gzip/bzip2/xz
  throughout.
\item \textbf{The lever is conditional volatility modeling, not the fractional
  form.} Against GARCH and FIGARCH the Mittag-Leffler kernel only ties or loses;
  it beats only EWMA. ``Long memory'' helps marginally (FIGARCH best by
  $\sim0.03\%$); the fractional clock is not special. The win over FLAC is robust
  to the volatility model.
\item \textbf{Scope is narrow and named.} The method loses to FLAC on speech by
  $\sim8\%$, and the comparison that would make ``beats FLAC'' into ``beats the
  ECG state of the art''---against wavelet / long-term-prediction ECG codecs---is
  not run here. The honest claim is: a volatility-adaptive arithmetic codec
  beats the general-purpose lossless standard on ECG by $3$--$5\%$.
\end{itemize}

\section*{Reproducibility}
All experiments are pure Python (NumPy/SciPy; \texttt{mpmath} for the
Mittag-Leffler function), with the real \texttt{flac} $1.4.3$ binary for the FLAC
baseline. Scripts: the volatility-model comparison (\texttt{lm\_vs\_garch.py});
the byte-level codec with exact round-trip and gzip/bzip2/xz baselines
(\texttt{bench.py}, with EWMA/GARCH/FIGARCH/Mittag-Leffler volatility models);
the parameter-fair FLAC comparison (\texttt{flac\_sweep.py}); and the
cross-recording tests (\texttt{ecg\_repeat.py}, \texttt{mit\_test.py}). ECG data:
the SciPy electrocardiogram and MIT-BIH record~100. The arithmetic coder is
round-trip verified lossless.

\begin{thebibliography}{9}
\bibitem{garch} T.~Bollerslev, \emph{Generalized autoregressive conditional
  heteroskedasticity}, J.\ Econometrics \textbf{31} (1986), 307--327.
\bibitem{figarch} R.~T.~Baillie, T.~Bollerslev, and H.~O.~Mikkelsen,
  \emph{Fractionally integrated generalized autoregressive conditional
  heteroskedasticity}, J.\ Econometrics \textbf{74} (1996), 3--30.
\bibitem{mem} R.~Fontes, \emph{The Fractional Memory Process: Well-Posedness,
  Spectral Decomposition, and Discrete Representation}, 2026.
\bibitem{GKMR} R.~Gorenflo, A.~A.~Kilbas, F.~Mainardi, and S.~V.~Rogosin,
  \emph{Mittag-Leffler Functions, Related Topics and Applications}, Springer,
  2014.
\bibitem{wnc} I.~H.~Witten, R.~M.~Neal, and J.~G.~Cleary, \emph{Arithmetic
  coding for data compression}, Comm.\ ACM \textbf{30} (1987), 520--540.
\bibitem{flac} J.~Coalson, \emph{FLAC---Free Lossless Audio Codec}, xiph.org.
\bibitem{mitbih} G.~B.~Moody and R.~G.~Mark, \emph{The impact of the MIT-BIH
  Arrhythmia Database}, IEEE Eng.\ Med.\ Biol.\ \textbf{20} (2001), 45--50.
\bibitem{physionet} A.~L.~Goldberger et al., \emph{PhysioBank, PhysioToolkit,
  and PhysioNet}, Circulation \textbf{101} (2000), e215--e220.
\bibitem{ecgcomp} M.~S.~Manikandan and S.~Dandapat, \emph{Wavelet-based
  electrocardiogram signal compression methods and their performances: a
  prospective review}, Biomed.\ Signal Process.\ Control \textbf{14} (2014),
  73--107.
\end{thebibliography}

\end{document}
